% GENERATED FILE. Edit canonical lessons, then run npm run generate:books.
% canonical-source-hash: fnv1a64:efc614f4ba9eae38
% canonical-lessons: ES-C16-imperfecto, ES-C16-comer-imperfecto, ES-C16-vivir-imperfecto

\chapter{The Imperfect --- What Used to Happen}
\label{ch:spanish-imperfecto}
\hlchaptermodality{80}

\begin{chapteropening}
\textbf{By the end of this chapter:} \emph{I can describe an ongoing or habitual past with a regular verb from any family.}

-aba and -ia, and -er and -ir agreeing once again.
\end{chapteropening}

\section[hablaba]{hablaba — let a past action stay open}

\label{lesson:ES-C16-imperfecto}



\begin{quote}
\textbf{Hablé español} presents speaking as a completed past event. A different past, the \textbf{imperfect}, can leave the action open: \textbf{Hablaba español} — “I was speaking Spanish” or “I used to speak Spanish.” The sentence does not say where the action ended.
\end{quote}

\begin{grammarlens}[title={Grammar Lens: three familiar persons}]
\noindent
\begin{tabularx}{\linewidth}{@{}>{\raggedright\arraybackslash}X>{\raggedright\arraybackslash}X@{}}
\toprule
\textbf{person} & \textbf{open or habitual past} \\
\midrule
yo & habl\textbf{aba} \\
tú & habl\textbf{abas} \\
él / ella / usted & habl\textbf{aba} \\
\bottomrule
\end{tabularx}

Say \textbf{hablaba · hablabas · hablaba}. The first and third forms are identical; the surrounding conversation or an explicit person tells you who is speaking. Plural persons wait.
\end{grammarlens}

\begin{cousinweb}
Spanish \textbf{-aba} continues the Latin imperfect family in \textbf{-ābam}. Historical sound change reshaped the complete Latin forms, but the \textbf{b} remains a useful audible bridge between \textbf{-ābam} and Spanish \textbf{-aba}.
\end{cousinweb}

\subsection*{Your turn}
Say these aloud:

\begin{itemize}
\raggedright
  \item “hablaba, hablabas, hablaba”
  \item “Hablé español. Hablaba español.”
  \item “-ābam and -aba share b”
\end{itemize}

\begin{tcolorbox}[breakable,colback=teal!4,colframe=teal!35!black,title={Before you move on}]
Give the three learned forms. (\emph{Hablaba, hablabas, hablaba.}) What can the imperfect do? (Leave a past action ongoing or habitual.) Which Latin family stands behind \textbf{-aba}? (\textbf{-ābam}.) Next: build the same gentle frame with \textbf{comer}.
\end{tcolorbox}
\section[comía]{comía — hear both vowels}

\label{lesson:ES-C16-comer-imperfecto}



\begin{quote}
\textbf{Comí} presents eating as completed. \textbf{Comía} leaves it open or habitual: “I was eating” or “I used to eat.” One accent separates the two forms in writing and in sound.
\end{quote}

\begin{grammarlens}[title={Grammar Lens: the singular -er row}]
\noindent
\begin{tabularx}{\linewidth}{@{}>{\raggedright\arraybackslash}X>{\raggedright\arraybackslash}X@{}}
\toprule
\textbf{person} & \textbf{open or habitual past} \\
\midrule
yo & com\textbf{ía} \\
tú & com\textbf{ías} \\
él / ella / usted & com\textbf{ía} \\
\bottomrule
\end{tabularx}

Say \textbf{comía · comías · comía}. As with \textbf{hablaba}, the first and third forms match. Plural persons wait.
\end{grammarlens}

\begin{grammarlens}[title={Grammar Lens: co-mí-a}]
The accent in \textbf{-ía} keeps \textbf{i} and \textbf{a} in separate syllables: \textbf{co-mí-a}, three beats. It appears in every form learned here. Compare final stress in \textbf{comí} with the middle stress and extra syllable in \textbf{comía}.
\end{grammarlens}

\begin{cousinweb}
The \textbf{-er} imperfect continues a Romance reshaping of Latin \textbf{-ēbam} material. Among the changes, an intervocalic \textbf{b} was lost; later Spanish spelling marks the stressed \textbf{í} that keeps \textbf{i-a} apart. This is a layered history, not a single instruction to delete one letter.
\end{cousinweb}

\subsection*{Your turn}
Say these aloud:

\begin{itemize}
\raggedright
  \item “comía, comías, comía”
  \item “Comí. Comía.”
  \item “co-mí-a”
\end{itemize}

\begin{tcolorbox}[breakable,colback=teal!4,colframe=teal!35!black,title={Before you move on}]
Give the three learned forms. (\emph{Comía, comías, comía.}) What does the accent do? (Keeps \textbf{i} and \textbf{a} apart.) Which Latin family lies behind the ending? (\textbf{-ēbam}, reshaped through Romance.) Next: test whether \textbf{vivir} can reuse this row.
\end{tcolorbox}
\section[vivía]{vivía — reuse the complete -er row}

\label{lesson:ES-C16-vivir-imperfecto}



\begin{quote}
An earlier chapter let \textbf{comer} and \textbf{vivir} share one completed-past row. They converge again here. Replace \textbf{com-} in \textbf{comía, comías, comía} with \textbf{viv-}.
\end{quote}

\begin{grammarlens}[title={Grammar Lens: one row reused}]
\noindent
\begin{tabularx}{\linewidth}{@{}>{\raggedright\arraybackslash}X>{\raggedright\arraybackslash}X>{\raggedright\arraybackslash}X@{}}
\toprule
\textbf{person} & \textbf{comer} & \textbf{vivir} \\
\midrule
yo & comía & vivía \\
tú & comías & vivías \\
él / ella / usted & comía & vivía \\
\bottomrule
\end{tabularx}

Say \textbf{vivía · vivías · vivía}. \textbf{-er} and \textbf{-ir} share \textbf{-ía, -ías, -ía}. The accent still keeps \textbf{i-a} apart. Plural persons wait.
\end{grammarlens}

\begin{grammarlens}[title={Grammar Lens: one known contrast}]
\begin{itemize}
\raggedright
  \item \textbf{Viví en Madrid.} — presents living there as a completed stretch.
  \item \textbf{Vivía en Madrid.} — puts us inside that past stretch or treats it as background.
\end{itemize}

Both are past. The speaker chooses how to frame the event.
\end{grammarlens}

\begin{cousinweb}
\textbf{Vivir} continues Latin \textbf{vīvere}. Its infinitive belongs to the Spanish \textbf{-ir} class, but the imperfect has converged with \textbf{-er}. That shared row is a reliable fact about these learned forms, not a claim that the two conjugations are identical everywhere.
\end{cousinweb}

\subsection*{Your turn}
Say these aloud:

\begin{itemize}
\raggedright
  \item “vivía, vivías, vivía”
  \item “comía and vivía share -ía”
  \item “Viví en Madrid. Vivía en Madrid.”
\end{itemize}

\begin{tcolorbox}[breakable,colback=teal!4,colframe=teal!35!black,title={Before you move on}]
Give the three \textbf{vivir} forms. (\emph{Vivía, vivías, vivía.}) Which class shares them? (\textbf{-er}.) What contrast can \textbf{viví / vivía} make? (Completed stretch versus an open or background past.) Next: learn \textbf{ver} before asking it to enter this tense.
\end{tcolorbox}
